\documentclass[11pt]{ctexart}
\usepackage[a4paper,margin=1in]{geometry}
\usepackage{graphicx}
\usepackage{xcolor}
\usepackage{hyperref}
\definecolor{refgray}{gray}{0.45}
\title{\textbf{市场最新观点汇总}\\[0.4em]{\large 全球资金正经历结构性再配置，AI与能源双引擎重塑美元定价，中国出口与消费的结构性升级成为核心主线。}}
\author{KC桌面 自动生成}
\date{260530}
\begin{document}
\maketitle
\section*{一页摘要}
\begin{itemize}
  \item 全球股票基金净流出70亿美元，但固定收益与货币市场基金持续流入，资金并非整体撤退，而是跨资产、跨区域、跨币种的系统性再配置。
  \item AI对存储的需求从GPU向CPU扩散，叠加HBM产能挤占与NAND供给纪律，存储TAM预计2027年达1.3万亿美元，市场尚未充分定价。
  \item 中国汽车出口结构质变，新能源渗透率4月首次突破55\%，插电混动占比从11\%跃升至39\%，成为海外市场最被低估的变量。
  \item 中国智能手机市场重心从出货量转向摄像头规格升级，2000万像素以上镜头渗透率达63\%，供应链价值分配逻辑重构。
  \item 美联储鹰派转向并非口头警告，核心PCE同比3.29\%且可能被低估，就业市场韧性掩盖真实工资压力，加息风险尚未被定价。
\end{itemize}
\section{宏观与利率}
\textbf{美联储鹰派转向反映通胀结构已变，就业韧性掩盖真实工资压力，加息风险尚未被市场充分定价。}

\begin{itemize}
  \item 4月核心PCE同比从3.24\%加速至3.29\%，远高于2\%目标，且美联储主席偏爱的截尾均值PCE可能低估实际通胀压力达48个基点。
  \item 5月非农就业预测新增11万人，但3个月移动均值升至13.7万，创2024年12月以来最强，劳动力市场韧性被技术性因素掩盖。
  \item 平均时薪4月放缓归因于日历效应，5月预计环比反弹至0.4\%，潜在工资增长仍高，市场对工资-通胀螺旋消退过于乐观。
  \item 市场仍在定价年内不降息也不加息的中性路径，但若通胀在及时框架内未回落，加息将成为政策工具箱中的现实选项。
\end{itemize}
{\small\color{refgray}\textbf{References:} [R010] 美联储的鹰派转向不是口头警告，而是通胀结构已变}

\section{AI与算力}
\textbf{AI需求从GPU向CPU扩散，存储市场供给约束强化，光纤光缆定价权迁移，人形机器人IPO潮开启供应链再定价。}

\begin{itemize}
  \item AI对存储的需求正从GPU向CPU扩散，CPU与GPU配比预计从2023年的5.4:1降至2028年的2.4:1，服务器级DDR/LPDDR5需求被大幅上修。
  \item 存储TAM在2027年预计达1.3万亿美元，长期协议合同销售占比上升将提升盈利稳定性，板块正从周期性商品向战略性资产转型。
  \item 日本光缆出口均价同比上涨77\%，光纤均价同比上涨44\%，价格涨幅远超数量涨幅，供给约束已从产能瓶颈转向定价权转移。
  \item 人形机器人行业进入资本化与规模化并行阶段，宇树科技等中国整机厂商提交IPO申请，供应链零部件需求将随产能扩张受益。
  \item Agentic AI对交易型垂直领域影响不均，电商等拥有不可替代供给的平台护城河更强，AI代理更可能放大而非取代其价值。
\end{itemize}
\begin{figure}[htbp]
\centering
\includegraphics[width=0.88\linewidth]{figures/fig_042.jpg}
\caption{Figure 1 - Figure 1: Aggregate market cap vs leading memory market revenue trend (1Q ahead)}
\end{figure}
\begin{figure}[htbp]
\centering
\includegraphics[width=0.88\linewidth]{figures/fig_043.jpg}
\caption{Figure 2 - Figure 2: Memory TAM vs peak mkt cap/TAM analysis US\textbackslash{}\$bn, peak mkt cap/TAM (x)}
\end{figure}
\begin{figure}[htbp]
\centering
\includegraphics[width=0.88\linewidth]{figures/fig_200.jpg}
\caption{Exhibit 5 - Exhibit 5: Average export price for optical fiber cable (¥1,000/kg)}
\end{figure}
\begin{figure}[htbp]
\centering
\includegraphics[width=0.88\linewidth]{figures/fig_201.jpg}
\caption{Exhibit 6 - Exhibit 6: Average export price for optical fiber (¥1,000/kg)}
\end{figure}
\begin{figure}[htbp]
\centering
\includegraphics[width=0.88\linewidth]{figures/fig_244.jpg}
\caption{Exhibit 5 - Exhibit 5: Unitree Revenue}
\end{figure}
{\small\color{refgray}\textbf{References:} [R005] 市场真正低估的不是需求，而是供给侧的再定价; [R014] 光纤光缆的真正拐点不在量，而在定价权的系统性迁移; [R016] 人形机器人IPO浪潮的真正含义：不是炒概念，而是供应链的再定价; [R011] 市场真正低估的不是AI需求，而是交易平台的护城河}

\section{能源与大宗}
\textbf{全球资源品供给硬约束以快于需求下滑的速度收紧，铝产能触及天花板，稀土定价分叉不可逆，中国定价权与西方安全溢价脱钩。}

\begin{itemize}
  \item 中国铝运行产能已达4530万吨，触及政府设定产能上限，供给弹性消失，即便利润高位也无法进一步扩张。
  \item 四月中国原铝产量390万吨同比增3.1\%，铝材出口环比增23\%同比增15\%，供给天花板与出口强劲共同支撑价格中枢。
  \item 中国稀土出口管制冲击集中于产业链纵深，5月钕铁硼磁材出口骤降77\%，对美出口暴跌92\%，西方供应链安全溢价正在重塑定价基准。
  \item 镨钕氧化物现货约103美元/公斤，低于MP Materials与Lynas的110美元价格下限，战略价格正成为西方供应新定价锚。
  \item 四月中国工业增加值同比增4.1\%低于预期，但铝是例外，供给端政策天花板与事故安全检查等硬约束成为比需求更关键的定价变量。
\end{itemize}
\begin{figure}[htbp]
\centering
\includegraphics[width=0.88\linewidth]{figures/fig_137.jpg}
\caption{Figure 2 - Figure 2: Historical NdPr Oxide price (inc \$13\textbackslash{}\%\$ VAT) and historical events}
\end{figure}
\begin{figure}[htbp]
\centering
\includegraphics[width=0.88\linewidth]{figures/fig_142.jpg}
\caption{Figure 8 - Figure 8: China Rare Earths Magnet Exports Monthly Heatmap Figure 9: Total China Rare Earth Magnets Export}
\end{figure}
\begin{figure}[htbp]
\centering
\includegraphics[width=0.88\linewidth]{figures/fig_145.jpg}
\caption{Figure 12 - Figure 12: Chinese rare earths magnet exports to US (t)}
\end{figure}
{\small\color{refgray}\textbf{References:} [R018] 市场真正低估的不是需求，而是供给侧的再定价; [R007] 稀土市场的结构性裂变：中国定价权与西方安全溢价正在脱钩}

\section{地产与消费}
\textbf{香港楼市脆弱点在于市场对内地购买力的结构性误读，中国智能手机与汽车出口的结构升级正在重塑供应链价值分配。}

\begin{itemize}
  \item 香港楼市所谓内地买家32\%交易金额占比中，超一半实际是已移居香港的居民，真正资金从内地过来的买家仅占交易额约15\%。
  \item 即便极端情景下所有内地购房需求消失，2026年全年住宅交易额仍能实现12\%同比增长，影响更集中于启德等特定新盘。
  \item 中国智能手机4月出货量同比增12\%环比增25\%，但二季度预计同比下滑6\%，存储芯片成本上涨压制需求。
  \item 智能手机摄像头规格升级趋势明确，平均摄像头数降至2.9颗创2020年新低，但2000万像素以上渗透率攀升至63\%创新高。
  \item 中国汽车出口新能源渗透率4月首次突破55\%，插电混动在新能源出口中占比从2024年初11\%跃升至39\%，成为海外市场黑马。
  \item 618美妆预售排名显示珀莱雅连续稳居第一，具备双阶段统治力，国际品牌首日表现更强，国货与中高端品牌后续持续发力。
\end{itemize}
\begin{figure}[htbp]
\centering
\includegraphics[width=0.88\linewidth]{figures/fig_126.jpg}
\caption{Figure 11 - Figure 1: 12m rolling no. of Mainland Chinese buyers for private residential units in HK}
\end{figure}
\begin{figure}[htbp]
\centering
\includegraphics[width=0.88\linewidth]{figures/fig_128.jpg}
\caption{Figure 3 - Figure 3: Percentage of Mainland Chinese buyers in private residential sales (by volume)}
\end{figure}
\begin{figure}[htbp]
\centering
\includegraphics[width=0.88\linewidth]{figures/fig_010.jpg}
\caption{Exhibit 12 - Exhibit 12: Huawei/Honor/Xiaomi/OPPO/Vivo/Transsion: 20MPx+ becomes the main contributor Cameras on smartphones launched by Huawei, Honor, Xiaomi, OPPO, Vivo, Transsion since 2018: \% of cameras in terms of pixels}
\end{figure}
\begin{figure}[htbp]
\centering
\includegraphics[width=0.88\linewidth]{figures/fig_023.jpg}
\caption{Figure 1 - Figure 1: China's PV exports: \% of BEV and PHEV of total, with April NEV penetration rate reaching 55\% units}
\end{figure}
\begin{figure}[htbp]
\centering
\includegraphics[width=0.88\linewidth]{figures/fig_024.jpg}
\caption{Figure 2 - Figure 2: PHEV \% of China's total PV NEV exports since 2024, with April reaching 39\%}
\end{figure}
{\small\color{refgray}\textbf{References:} [R006] 香港楼市真正的脆弱点不是资本管制，而是市场如何重新定价“内地购买力”; [R002] 中国智能手机市场真正被低估的，不是4月出货量反弹，而是摄像头规格升级的结构性定价权; [R004] 中国汽车出口的真正转折点：不是规模，而是结构; [R017] 618预售排名真正透露的信号：不是消费降级，而是品牌力的结构性重排}

\section{区域市场}
\textbf{美元中期支撑强于市场预期，AI繁荣与能源供给冲击双引擎重塑定价锚，日本半导体材料与印刷巨头估值错配提供窗口。}

\begin{itemize}
  \item 美元近期下跌是战术性而非结构性，AI繁荣吸引外资回流美元资产，能源供给冲击持续推动增长分化，美元中期支撑强于预期。
  \item 韩国资金持续涌入美国股票，贸易顺差未回流国内经济，经常账户改善不再支撑韩元，传统外汇分析框架被资本账户金融化颠覆。
  \item 日本TOPPAN和DNP被市场错误归类为传统印刷商，FC-BGA业务营收CAGR达35\%，当前P/B不到1倍提供估值错配窗口。
  \item 日本MLCC四月出货金额同比增28\%，ASP同比涨15.9\%环比涨2.6\%，AI服务器需求推动定价脱离传统周期特征走向结构性提升。
  \item NAND设备支出拐点正在萌芽，Kioxia资本开支指引4500亿日元，但短期证据有限，市场需耐心等待2027年加速。
\end{itemize}
\begin{figure}[htbp]
\centering
\includegraphics[width=0.88\linewidth]{figures/fig_191.jpg}
\caption{Exhibit 3 - Exhibit 3: FC-BGA sales growth: FY24=100 FC-BGA revenue}
\end{figure}
\begin{figure}[htbp]
\centering
\includegraphics[width=0.88\linewidth]{figures/fig_192.jpg}
\caption{Exhibit 4 - Exhibit 4: FC-BGA operating profit growth: FY24=100 FC-BGA OP}
\end{figure}
\begin{figure}[htbp]
\centering
\includegraphics[width=0.88\linewidth]{figures/fig_001.jpg}
\caption{Figure 1 - Takayuki Naito \$\textasciicircum{}\{AC\}\$}
\end{figure}
{\small\color{refgray}\textbf{References:} [R012] 美元并未真正走弱，市场低估的是“双引擎”对美元定价权的重塑; [R003] 市场真正低估的不是需求，而是资本流动的结构性再定价; [R013] 日本印刷巨头的真正价值不在印刷，而在半导体封装材料的定价权; [R001] 日本MLCC四月数据真正揭示的，不是需求回暖，而是ASP的结构性重塑; [R019] 市场真正低估的不是需求，而是NAND和Intel的拐点何时到来}

\section{风险偏好与资金流向}
\textbf{全球资金进行系统性再配置，2026年IPO创纪录但企业股权需求超过供给，中国生物科技从证明全球化转向锁定价值留存。}

\begin{itemize}
  \item 截至5月27日当周全球股票基金净流出70亿美元，但固定收益基金净流入240亿美元，货币市场基金资产增220亿美元。
  \item 2026年IPO规模预计2250亿美元创历史新高，但企业股权总供给仅占罗素3000总市值1.0\%，低于1995年以来1.5\%均值。
  \item 企业自身股权需求将超过供给，回购公告已达8600亿美元，全年预计1.3万亿美元，净回购收益率虽降至1.9\%仍是重要需求来源。
  \item 中国生物科技下一阶段重估奖励的不是交易数量，而是风险承担能力与资产稀缺性共同决定的价值留存率。
  \item 中国生物科技对外授权交易与并购交易相比，原研方留存的风险调整后净现值存在约51\%折价，NewCo或共同开发结构可收窄至35-41\%。
  \item 字节跳动系App用户时长占比突破38.3\%首次超越腾讯，在线视频、音乐、阅读等多赛道同时发起系统性挑战。
\end{itemize}
\begin{figure}[htbp]
\centering
\includegraphics[width=0.88\linewidth]{figures/fig_202.jpg}
\caption{Exhibit 1 - Exhibit 1: 40 US companies have gone public so far this year US deals greater than \textbackslash{}\$25 million in value}
\end{figure}
\begin{figure}[htbp]
\centering
\includegraphics[width=0.88\linewidth]{figures/fig_205.jpg}
\caption{Exhibit 35 - Exhibit 4: We forecast \textbackslash{}\$225 billion of IPO gross proceeds in 2026}
\end{figure}
\begin{figure}[htbp]
\centering
\includegraphics[width=0.88\linewidth]{figures/fig_206.jpg}
\caption{Exhibit 5 - Exhibit 5: We expect equity issuance will equal 1\% of Russell 3000 market cap this year}
\end{figure}
\begin{figure}[htbp]
\centering
\includegraphics[width=0.88\linewidth]{figures/fig_150.jpg}
\caption{Exhibit 3 - Exhibit 3: Innovation Dawn 2.0 positioning: risk-bearing capacity vs. novel asset scarcity}
\end{figure}
\begin{figure}[htbp]
\centering
\includegraphics[width=0.88\linewidth]{figures/fig_151.jpg}
\caption{Exhibit 4 - Exhibit 4: Out-licensing to M\&A NPV discount (US\textbackslash{}\$mn)}
\end{figure}
{\small\color{refgray}\textbf{References:} [R003] 市场真正低估的不是需求，而是资本流动的结构性再定价; [R015] 2026年IPO创纪录，但市场真正该担心的不是供给过剩，而是企业需求的结构性变化; [R008] 中国生物科技的下一个重估：从“证明全球化”到“锁定价值留存”; [R009] 字节跳动正在系统性重写中国互联网的竞争规则}

\section*{结语}
综合来看，市场主线正从需求侧叙事转向供给侧结构性变化与资本流动再定价，AI与能源双引擎重塑全球资产定价框架，中国出口与消费的结构升级成为不可忽视的alpha来源。
\vfill
{\small\color{refgray}Personal reading notes and learning share only. Not investment advice.}
\end{document}
